\documentclass[11pt,a4paper]{article}

\usepackage[T1]{fontenc}
\usepackage[utf8]{inputenc}
\usepackage{times}
\usepackage[spanish,es-noshorthands,es-tabla]{babel}
\usepackage{csquotes}
\usepackage{amsmath,amssymb}
\usepackage{booktabs}
\usepackage{tabularx}
\usepackage{array}
\usepackage{enumitem}
\usepackage{setspace}
\usepackage{xcolor}
\usepackage{graphicx}
\usepackage{tikz}
\usetikzlibrary{positioning}
\usepackage[margin=2.54cm]{geometry}
\usepackage{titlesec}
\usepackage{caption}
\usepackage{fancyhdr}
\usepackage[backend=biber,style=apa]{biblatex}
\DeclareLanguageMapping{spanish}{spanish-apa}
\addbibresource{referencias.bib}
\usepackage[hidelinks]{hyperref}

\onehalfspacing
\setlength{\parindent}{1.27cm}
\setlength{\parskip}{0pt}

% ---- Encabezados estilo APA 7 ----
\titleformat{\section}{\normalfont\normalsize\bfseries\centering}{\thesection.}{0.5em}{}
\titleformat{\subsection}{\normalfont\normalsize\bfseries}{\thesubsection}{0.5em}{}
\titleformat{\subsubsection}{\normalfont\normalsize\bfseries\itshape}{\thesubsubsection}{0.5em}{}
\titlespacing*{\section}{0pt}{12pt}{6pt}
\titlespacing*{\subsection}{0pt}{10pt}{4pt}
\titlespacing*{\subsubsection}{0pt}{8pt}{3pt}

% ---- Tablas y figuras estilo APA ----
\captionsetup[table]{labelfont=bf,textfont=it,labelsep=newline,justification=raggedright,singlelinecheck=false,position=top}
\captionsetup[figure]{labelfont=bf,textfont=it,labelsep=newline,justification=raggedright,singlelinecheck=false,position=top}

% ---- Número de página arriba a la derecha (APA) ----
\pagestyle{fancy}\fancyhf{}\fancyhead[R]{\thepage}\renewcommand{\headrulewidth}{0pt}

\newcommand{\tabfont}{\small\setstretch{1}}

\begin{document}

% =====================================================================
%                              PORTADA
% =====================================================================
\begin{center}
{\large\bfseries Optimización basada en inteligencia de enjambre (PSO y Bat Algorithm) para la selección de características en la detección de intrusiones sobre tráfico real de un enlace troncal transpacífico}\\[10pt]
Medina Fuentes, Jhojan Sebastian\\
Facultad de Ingeniería Eléctrica y Electrónica, Universidad Nacional de Ingeniería\\
Inteligencia Artificial II, Prof.\ Yury Oscar Tello Canchapoma\\
2026
\end{center}

\subsection*{Resumen}
Los sistemas de detección de intrusiones (IDS) basados en flujos de red procesan decenas de atributos, muchos de ellos redundantes, lo que eleva el costo computacional y puede aumentar las falsas alarmas. Este estudio evalúa si la selección de características con Optimización por Enjambre de Partículas (PSO) y con el Algoritmo del Murciélago (BA), ambos binarios y en esquema \emph{wrapper}, mejora el rendimiento de detección frente a un clasificador Random Forest que usa todas las características (línea base de 50 atributos) y un control aleatorio (\emph{Random-k21}). Se emplea tráfico real y crudo del archivo MAWI (enlace troncal transpacífico del proyecto WIDE, punto de muestreo F) con etiquetas de MAWILab v1.1. Los flujos se extraen directamente de los archivos PCAP de 3 capturas de 2021 ($N = 3{,}000{,}000$ flujos muestreados con 22,82\,\% de tráfico de ataque) y se comparan los esquemas con diseño pareado ($N = 30$ réplicas independientes, semillas 1 a 30). Se miden la eficacia ($F_1$ macro, recall de la clase ataque y accuracy), las falsas alarmas (tasa de falsos positivos, TFP, y de falsos descubrimientos, TFD) y la eficiencia computacional (reducción dimensional y tiempos de entrenamiento e inferencia). Las diferencias se contrastan con las pruebas de Friedman y Wilcoxon con ajuste post-hoc de Holm ($\alpha = 0{,}05$) y tamaño del efecto $r_b$. PSO redujo los atributos a 18,5 de media (reducción del 63,1\,\%) pero no mejoró significativamente el $F_1$ macro frente a la línea base (0,734 vs. 0,728; $p_{\text{Holm}} = 1{,}0000$), sufriendo una caída significativa en el recall de la clase ataque de 6,39 puntos porcentuales (8,98\,\% relativo; 0,648 vs. 0,712; $p_{\text{Holm}} < 0{,}001$). BA resultó significativamente inferior en $F_1$ y recall. Respecto a falsas alarmas, PSO redujo la TFP de 0,192 a 0,152 ($-21{,}2\,\%$), sin alcanzar significancia estadística tras la corrección de Holm ($p_{\text{Holm}} = 0{,}5264$). El beneficio principal fue computacional: PSO disminuyó el tiempo de entrenamiento en 36,07\,\% ($p = 0{,}00002$) y la latencia de inferencia en 10,75\,\% ($p = 0{,}00015$). Sin embargo, el control aleatorio \emph{Random-k21} redujo el tiempo de entrenamiento en 36,75\,\% sin diferir de PSO ($p = 0{,}0919$), demostrando que el ahorro temporal obedece a la reducción dimensional per se y no a la metaheurística. Se concluye que en tráfico troncal real no se comprueba la hipótesis de superioridad en eficacia reportada usualmente en bancos sintéticos.

\medskip
\noindent\textbf{Palabras clave:} detección de intrusiones, selección de características, optimización por enjambre de partículas, algoritmo del murciélago, tráfico real, MAWILab.

\subsection*{Ideas destacadas}
\begin{itemize}[leftmargin=1.5em,itemsep=0pt,topsep=2pt]
  \item PSO y BA se comparan con una línea base y un control aleatorio bajo condiciones idénticas sobre tráfico troncal real.
  \item Los atributos de flujo se extraen directamente de PCAP originales, sin conjuntos preprocesados de laboratorio.
  \item La selección con PSO reduce 63\,\% de atributos y 36\,\% de tiempo de cómputo, pero no supera a la línea base en $F_1$ macro ni en TFP tras la corrección estadística estricta.
\end{itemize}

% =====================================================================
\section{Introducción}

% --- Contexto ---
Los sistemas de detección de intrusiones en red (IDS) analizan flujos de tráfico descritos por decenas de atributos estadísticos. La selección de características es un componente habitual en el diseño de IDS eficientes \parencite{zhou2020}, y los algoritmos de enjambre y evolutivos se han estudiado de forma amplia como herramientas de búsqueda para ese fin \parencite{thakkar2020}.

% --- Problemática de la VD ---
El rendimiento de detección de un IDS no depende solo de acertar sobre el tráfico de ataque. También depende de cuántas falsas alarmas genera y de cuánto cuesta, en tiempo y en memoria, entrenar el modelo y clasificar cada flujo. Cuando los atributos son numerosos, redundantes o ruidosos, aumenta el costo computacional y puede deteriorarse la capacidad de distinguir el ataque del tráfico benigno \parencite{zhou2020,thakkar2020}. En entornos de enlace troncal transpacífico con flujos continuos a escala de gigabits por segundo, un incremento marginal en la tasa de falsos positivos satura la capacidad operativa de los analistas de seguridad, donde una tasa aparentemente baja del 1\,\% representa decenas de miles de alertas erróneas al día. El problema que aborda este trabajo es, por tanto, el deficiente rendimiento de detección del tráfico de red en un IDS de enlace troncal transpacífico, asociado a la alta dimensionalidad de los atributos de flujo.

% --- Estado de la técnica ---
Entre las estrategias para reducir la dimensionalidad, el método \emph{wrapper} evalúa cada subconjunto de atributos entrenando un clasificador, y se combina de forma natural con metaheurísticas de búsqueda. La literatura reporta propuestas basadas en PSO \parencite{srinoy2007,louk2022,safa2024}, en variantes del algoritmo del murciélago \parencite{ghanem2022,saheed2024}, en algoritmos genéticos \parencite{halim2021}, en optimizadores de inspiración natural distinta o híbridos \parencite{vijayanand2020,alzubi2022,hassan2022,varzaneh2024} y en optimización combinatoria \parencite{nazir2021}.

% --- Vacío y aporte ---
La validez práctica de estas propuestas depende de los datos con que se evalúan. La literatura dominante en selección de características bio-inspirada evalúa sus métodos en bancos de laboratorio o sintéticos clásicos (KDD Cup 99, NSL-KDD, UNSW-NB15 y CIC-IDS2017) \parencite{thakkar2020,alzubi2022,hassan2022,safa2024}. Dichos repositorios presentan distribuciones predecibles y ataques sintéticos con huellas homogéneas, los cuales no reflejan la variabilidad estadística, el ruido de fondo ni el desplazamiento distributivo (\emph{covariate shift}) propio de enlaces troncales operativos. El archivo MAWI reúne trazas de tráfico real del enlace troncal transpacífico WIDE \parencite{mawiarchive}, con etiquetas de anomalías obtenidas por consenso automatizado de detectores \parencite{fontugne2010}. Hasta donde alcanza la revisión realizada para este trabajo, no se identificó una comparación empírica y rigurosa de PSO y BA para selección de características ejecutada directamente sobre flujos extraídos de trazas crudas de MAWI bajo protocolos pareados con $N \ge 30$ réplicas y corrección de comparaciones múltiples. El aporte de este estudio es evaluar esa comparación bajo condiciones controladas, con un conjunto de flujos construido de forma propia y un diseño que permite contrastes estadísticos pareados.

% --- Propósito, objetivos, motivación, importancia ---
El \textbf{propósito} es determinar en qué medida la selección de características basada en inteligencia de enjambre (PSO y BA) mejora el rendimiento de detección del tráfico de red en dicho IDS, frente a un clasificador que usa todas las características. Los \textbf{objetivos específicos} son evaluar su efecto sobre (i) la eficacia de detección ($F_1$ macro, recall de la clase ataque y accuracy), (ii) las falsas alarmas (tasa de falsos positivos y de falsos descubrimientos) y (iii) la eficiencia computacional (reducción de atributos y de los tiempos de entrenamiento e inferencia).

La \textbf{motivación} es que un IDS que opere sobre un enlace troncal debe procesar grandes volúmenes de tráfico, de modo que reducir atributos sin perder capacidad de detección tiene un efecto directo en su costo de operación. La \textbf{importancia} del estudio radica en que sus conclusiones se apoyan en tráfico real y crudo y no en datos generados en laboratorio, aportando evidencia empírica que desmitifica la transferencia directa de hiper-rendimientos desde benchmarks sintéticos hacia infraestructuras de telecomunicaciones reales. El artículo se organiza así: la Sección 2 revisa la literatura, la 3 describe la metodología, la 4 detalla los algoritmos, y las Secciones 5 a 7 presentan los resultados, la discusión y las conclusiones.

% =====================================================================
\section{Revisión de la literatura}

\textcite{thakkar2020} revisan exhaustivamente el papel de los algoritmos de enjambre y evolutivos en los sistemas de detección de intrusiones, identificando que la gran mayoría de publicaciones recurre a conjuntos de datos estáticos y sintéticos (KDD Cup 99, NSL-KDD), omitiendo el impacto del tráfico de alta velocidad. Dentro de esa línea, \textcite{safa2024} combinan PSO con Random Forest para seleccionar atributos relevantes en un IDS sobre conjuntos públicos sintéticos, alcanzando mejoras de hasta 3,4\,\% en precisión, aunque con tamaños de muestra limitados ($N \le 5$) y sin validar la latencia de inferencia por flujo. Por su parte, \textcite{louk2022} proponen una selección de características guiada por PSO junto con un ensamble híbrido sobre CIC-IDS2017, reportando alta eficacia sin incorporar pruebas de hipótesis pareadas ni contrastar algoritmos de enjambre concurrentes.

Para el algoritmo del murciélago, \textcite{ghanem2022} presentan un BA binario multiobjetivo para selección de características y optimización de redes neuronales sobre NSL-KDD y UNSW-NB15, logrando exactitudes superiores al 99\,\% que no han sido contrastadas en enlaces transpacíficos. Asimismo, \textcite{saheed2024} acoplan BA con sistemas de numeración de residuos para filtrar atributos en NSL-KDD, ratificando la necesidad de examinar metaheurísticas bio-inspiradas en escenarios con ruido de fondo no supervisado. La Tabla~\ref{tab:sintesis} resume los cinco estudios clave que evidencian el vacío metodológico.

\begin{table}[ht]
\tabfont
\caption{Síntesis de los estudios que evidencian el vacío}
\label{tab:sintesis}
\begin{tabularx}{\textwidth}{@{}p{2.6cm}p{3.2cm}p{2.8cm}p{2.3cm}X@{}}
\toprule
Estudio & Técnica de selección & Clasificador & Datos & Limitación frente a este trabajo \\
\midrule
\textcite{thakkar2020} & Revisión de enjambre y evolutivos & Varios clasificadores & Benchmarks sintéticos & Revisión teórica sin experimentación en tráfico troncal real \\
\textcite{safa2024} & PSO & Random Forest & Benchmarks públicos & Réplicas reducidas ($N \le 5$); sin benchmark de latencia por flujo \\
\textcite{louk2022} & PSO & Conjunto híbrido & CIC-IDS2017 & Tráfico de laboratorio; sin contrastes post-hoc corregidos \\
\textcite{ghanem2022} & BA binario multiobjetivo & Red neuronal & NSL-KDD, UNSW-NB15 & No compara con PSO bajo el mismo modelo en enlace troncal \\
\textcite{saheed2024} & BA híbrido (RNS) & Random Forest & NSL-KDD & Datos sintéticos sin deriva temporal ni validación LOO \\
\bottomrule
\end{tabularx}
\end{table}

% =====================================================================
\section{Metodología}

\subsection{Diseño y justificación del método}
El estudio es cuantitativo, aplicado y de nivel correlacional-causal, con un diseño cuasi-experimental transversal. Se eligió un experimento computacional porque permite controlar los factores que podrían confundir la comparación: los cuatro brazos evaluados (línea base de 50 atributos, control aleatorio \emph{Random-k21}, PSO más Random Forest y BA más Random Forest) usan exactamente el mismo clasificador (\emph{Random Forest}, 100 estimadores, profundidad 15, ponderación balanceada), la misma partición estratificada de datos y las mismas semillas estocásticas (1 a 30). Las hipótesis formales de trabajo formuladas en la matriz de consistencia son:
\begin{itemize}[leftmargin=1.5em,itemsep=1pt]
  \item[\textbf{H1:}] La selección de características mediante PSO mejora significativamente el $F_1$ macro respecto a la línea base.
  \item[\textbf{H2:}] La selección con PSO reduce significativamente la tasa de falsos positivos (TFP).
  \item[\textbf{H3:}] La reducción de dimensionalidad reduce el tiempo de entrenamiento e inferencia en al menos un 20\,\% (umbral de la matriz de operacionalización; se analiza también la exigencia al 25\,\%).
  \item[\textbf{H4:}] PSO y BA producen desempeños cualitativamente diferenciables en métricas clave.
\end{itemize}

\subsection{Dataset}
Se emplean capturas PCAP crudas del archivo MAWI \parencite{mawiarchive}, punto de muestreo F (enlace de tránsito de 150 Mbps de WIDE hacia su proveedor transpacífico en Japón), con etiquetas de MAWILab v1.1 \parencite{fontugne2010}. Cada archivo corresponde a 15 minutos de tráfico ininterrumpido conservando las cabeceras TCP/IP anonimizadas. Se seleccionaron 3 días de 2021 pertenecientes a distintos meses y franjas horarias consistentes (14:00 JST), antes del cese de actualización de MAWILab (Tabla~\ref{tab:capturas}).

\begin{table}[ht]
\tabfont
\caption{Capturas empleadas (MAWI, punto de muestreo F, 2021)}
\label{tab:capturas}
\begin{tabular}{@{}lrrrr@{}}
\toprule
Fecha & PCAP (GB) & Flujos extraídos & Ataque (\%) & Benigno (\%) \\
\midrule
2021-02-10 & 2,65 & 1.000.000 & 22,82 & 77,18 \\
2021-06-16 & 1,69 & 1.500.000 & 22,82 & 77,18 \\
2021-10-13 & 1,81 &   500.000 & 22,82 & 77,18 \\
\midrule
\textbf{Total} & 6,15 & 3.000.000 & 22,82 & 77,18 \\
\bottomrule
\end{tabular}

\smallskip
\footnotesize\emph{Nota.} Tamaño de los archivos comprimidos (.pcap.gz). Las muestras conservan la proporción natural de anomalías validadas.
\end{table}

\subsection{Extracción de flujos y etiquetado}
Los atributos se extraen de los PCAP con un extractor propio en Python optimizado con \texttt{dpkt} 1.9.8. Un flujo se define como la secuencia bidireccional de paquetes que comparten la tupla de 5 elementos (\emph{IP origen}, \emph{IP destino}, \emph{puerto origen}, \emph{puerto destino}, \emph{protocolo}), con expiración tras 64 segundos de inactividad o un máximo de 120 segundos de duración, así como ante el cierre TCP por banderas FIN o RST. Como las trazas crudas truncan la carga útil, las longitudes se obtienen del encabezado IP. Se calculan 50 atributos numéricos que abarcan duración, paquetes y bytes por sentido, longitudes de paquete, tiempos entre arribos (IAT), conteo de banderas TCP y razones de direccionalidad.

El etiquetado se realiza a partir de los archivos XML de MAWILab v1.1, cruzando las anomalías categorizadas como \emph{anomalous} y \emph{suspicious} (definidas formalmente como clase ataque = 1). Los flujos clasificados exclusivamente como \emph{notice} se descartaron por corresponder a eventos aislados no ratificados por consenso de detectores. Cada anomalía se proyecta sobre ambos sentidos del flujo considerando comodines y marcas temporales de inicio y fin.

\subsection{Preprocesamiento y partición}
Los campos identificadores (\emph{IPs}, puertos, marcas temporales y claves de flujo) se excluyen del conjunto de aprendizaje para evitar atajos espurios (\emph{data leakage}). Los valores infinitos y nulos generados por divisiones entre cero en ráfagas cortas se imputan con 0,0. Se verificó que ninguno de los 50 atributos resultara constante. El conjunto consolidado de $3{,}000{,}000$ flujos se divide de manera estratificada en 70\,\% para entrenamiento ($2{,}100{,}000$ flujos) y 30\,\% para prueba ($900{,}000$ flujos) fijando la semilla global en 42. Para evaluar la aptitud durante las iteraciones metaheurísticas se utiliza una submuestra estratificada de 10.000 instancias del conjunto de entrenamiento. Adicionalmente, se construyeron particiones \emph{Leave-One-Capture-Out} (LOO) para validar la resistencia de los subconjuntos seleccionados frente al \emph{covariate shift} temporal entre días.

\subsection{Diseño experimental y métricas}
El experimento comprende 30 ejecuciones independientes ($N = 30$) con semillas $k \in [1, 30]$ compartidas por todos los brazos. En cada réplica $k$: (1) se evalúa la línea base con 50 atributos; (2) se evalúa \emph{Random-k21} con 21 atributos al azar; (3) se ejecutan PSO y BA binarios para seleccionar atributos; (4) se reentrena Random Forest en el 100\,\% del conjunto de entrenamiento con el subconjunto obtenido; y (5) se infiere sobre los $900{,}000$ flujos de prueba. Con la matriz de confusión ($VP$, $VN$, $FP$, $FN$, donde la clase positiva es ataque) se computan:
\begin{align}
F1_{\text{macro}} &= \frac{1}{2}\sum_{c=0}^{1}\frac{2P_cR_c}{P_c+R_c}, & \text{Recall}_{\text{ataque}}&=\frac{VP}{VP+FN}, & \text{Accuracy}&=\frac{VP+VN}{N},\\
TFP &= \frac{FP}{FP+VN}, & TFD &= \frac{FP}{FP+VP}, & RD &= \frac{50-N_{\text{sel}}}{50}\times100,
\end{align}
junto con la reducción porcentual del tiempo de entrenamiento ($RTE$) y de la latencia de inferencia ($RTI$). Debe aclararse que la métrica de sensibilidad evaluada corresponde estrictamente al \emph{Recall de la clase ataque} (tasa de verdaderos positivos sobre flujos maliciosos, $VP/(VP+FN)$) y no a la media macro de recall de ambas clases.

\subsection{Análisis estadístico}
Siguiendo a \textcite{demsar2006}, se aplica la prueba no paramétrica de Friedman sobre las 30 réplicas emparejadas. Si $\chi_F^2$ resulta significativo ($\alpha = 0{,}05$), se procede con la prueba de Wilcoxon pareada con corrección escalonada de Holm y tamaño del efecto $r_b$ (correlación biserial de rangos). Para diferencias pareadas $d = X_A - X_B$, se reporta $r_b$ estandarizado; en tiempos de ejecución, un valor negativo de $d$ indica que el brazo $A$ es más rápido (requiere menor tiempo).

% =====================================================================
\section{Desarrollo}

\subsection{Arquitectura general}
La Figura~\ref{fig:arq} describe el flujo integral del sistema, desde la ingesta de paquetes brutos hasta los contrastes estadísticos automatizados.

\begin{figure}[ht]
\centering
\caption{Arquitectura del sistema experimental}
\label{fig:arq}
\begin{tikzpicture}[box/.style={draw,rounded corners,align=center,minimum height=0.95cm,minimum width=2.5cm,font=\footnotesize},node distance=0.55cm and 0.6cm,>=stealth]
\node[box] (a) {PCAP crudos\\MAWI (F, 2021)};
\node[box,right=of a] (b) {Extracción\\de flujos};
\node[box,right=of b] (c) {Etiquetado\\MAWILab};
\node[box,right=of c] (d) {Preprocesamiento\\y partición 70/30};
\node[box,below=0.9cm of d] (e) {Búsqueda wrapper\\PSO / BA};
\node[box,left=of e] (f) {Random Forest\\(aptitud, CV)};
\node[box,left=of f] (g) {Evaluación final\\(prueba 30\,\%)};
\draw[->] (a)--(b); \draw[->] (b)--(c); \draw[->] (c)--(d); \draw[->] (d)--(e);
\draw[<->] (e)--(f); \draw[->] (f)--(g);
\end{tikzpicture}
\end{figure}

\subsection{Clasificador y metaheurísticas}
Se emplea Random Forest \parencite{breiman2001} con 100 estimadores (\texttt{max\_depth = 15}, \texttt{class\_weight = 'balanced'}). La aptitud a minimizar es:
\begin{equation}
f(\mathbf{x}) = 0{,}95\,\bigl(1 - F1_{\text{macro}}^{\text{CV}}(\mathbf{x})\bigr) + 0{,}05\,\frac{|\mathbf{x}|}{50}.
\label{eq:fitness}
\end{equation}
PSO binario opera con $\omega = 0{,}70$, $c_1 = c_2 = 1{,}50$, $v_{\max} = 6{,}0$ y sigmoide \parencite{kennedy1995}. BA binario opera con $f \in [0, 2]$, $A_0 = 1{,}0$, $r_0 = 0{,}5$ y $\alpha = \gamma = 0{,}90$ \parencite{yang2010}. Ambos corren con 10 individuos por 15 iteraciones sobre una submuestra de 10.000 flujos.

% =====================================================================
\section{Resultados}

\subsection{Datos finales}
Se procesaron $6{,}15\text{ GB}$ brutos de tráfico troncal MAWI de 2021 correspondientes a $3{,}000{,}000$ flujos consolidados. El 22,82\,\% correspondió a anomalías validadas y el 77,18\,\% a tráfico benigno. La exactitud trivial fue de 77,18\,\%, con un $F_1$ macro base de 0,4356.

\subsection{Eficacia y falsas alarmas}
La Tabla~\ref{tab:efi} presenta el desempeño de los cuatro brazos sobre los $900{,}000$ flujos del conjunto de prueba en las 30 repeticiones.

\begin{table}[ht]
\tabfont
\caption{Eficacia y falsas alarmas por brazo (media $\pm$ desviación estándar; $N = 30$)}
\label{tab:efi}
\begin{tabular}{@{}lccccc@{}}
\toprule
Brazo & $F1$ macro & Recall (clase ataque) & Accuracy & TFP & TFD \\
\midrule
Línea base (50 atributos) & $0{,}7280 \pm 0{,}0001$ & $0{,}7117 \pm 0{,}0004$ & $0{,}7858 \pm 0{,}0002$ & $0{,}1924 \pm 0{,}0003$ & $0{,}4776 \pm 0{,}0003$ \\
Random-k21 (control)      & $0{,}7161 \pm 0{,}0373$ & $0{,}7071 \pm 0{,}1036$ & $0{,}7739 \pm 0{,}0553$ & $0{,}2063 \pm 0{,}1020$ & $0{,}4716 \pm 0{,}0873$ \\
PSO + Random Forest       & $\mathbf{0{,}7342 \pm 0{,}0155}$ & $0{,}6478 \pm 0{,}0764$ & $\mathbf{0{,}8026 \pm 0{,}0281}$ & $\mathbf{0{,}1516 \pm 0{,}0588}$ & $\mathbf{0{,}4202 \pm 0{,}0923}$ \\
BA + Random Forest        & $0{,}7244 \pm 0{,}0016$ & $0{,}6993 \pm 0{,}0175$ & $0{,}7840 \pm 0{,}0029$ & $0{,}1910 \pm 0{,}0086$ & $0{,}4801 \pm 0{,}0058$ \\
Clasificador trivial      & 0,4356 & 0,0000 & 0,7718 & 0,0000 & — \\
\bottomrule
\end{tabular}
\end{table}

\subsection{Subconjuntos seleccionados y benchmark formal de tiempos}
La Tabla~\ref{tab:efic} resume la reducción dimensional y el benchmark formal de tiempos ($N=30$). PSO seleccionó en promedio 18,5 atributos (reducción dimensional del 63,1\,\%) y BA 24,2 atributos (51,5\,\%).

\begin{table}[ht]
\tabfont
\caption{Eficiencia computacional por brazo ($N = 30$; media $\pm$ desviación estándar)}
\label{tab:efic}
\begin{tabular}{@{}lccccc@{}}
\toprule
Brazo & Atributos ($N_{\text{sel}}$) & RD (\%) & $t_{\text{entrena}}$ (s) & RTE (\%) & RTI (\%) \\
\midrule
Línea base (50 atrib.) & $50{,}0 \pm 0{,}0$ & 0,0 & $26{,}45 \pm 15{,}21$ & — & — \\
Random-k21 (control)   & $21{,}0 \pm 0{,}0$ & 58,0 & $16{,}73 \pm 8{,}67$ & +36,75 & +8,23 \\
PSO + Random Forest    & $18{,}5 \pm 2{,}6$ & 63,1 & $16{,}91 \pm 8{,}96$ & \textbf{+36,07} & \textbf{+10,75} \\
BA + Random Forest     & $24{,}2 \pm 3{,}0$ & 51,5 & $19{,}45 \pm 10{,}71$ & +26,47 & +3,69 \\
\bottomrule
\end{tabular}

\smallskip
\footnotesize\emph{Nota.} Latencia de inferencia: Línea base $= 2{,}256\,\mu\text{s/flujo}$; Random-k21 $= 2{,}071\,\mu\text{s/flujo}$; PSO $= 2{,}014\,\mu\text{s/flujo}$; BA $= 2{,}173\,\mu\text{s/flujo}$. Tiempo medio de búsqueda metaheurística: PSO $= 307{,}8\,\text{s}$; BA $= 42{,}3\,\text{s}$.
\end{table}

\noindent\textbf{Análisis de amortización computacional:}
El tiempo invertido en la búsqueda metaheurística debe recuperarse mediante el ahorro acumulado durante los sucesivos reentrenamientos del modelo:
\begin{itemize}[leftmargin=1.5em,itemsep=1pt]
  \item \textbf{PSO:} Frente a la línea base, ahorra $26{,}45 - 16{,}91 = 9{,}54\text{ s}$ por entrenamiento. Requiere $\lceil 307{,}8 / 9{,}54 \rceil = \mathbf{33\text{ ciclos}}$ de reentrenamiento para amortizar los $307{,}8\text{ s}$ de búsqueda.
  \item \textbf{BA:} Frente a la línea base, ahorra $26{,}45 - 19{,}45 = 7{,}00\text{ s}$ por entrenamiento. La función techo determina que $\lceil 42{,}3 / 7{,}00 \rceil = \lceil 6{,}04 \rceil = \mathbf{7\text{ ciclos}}$ son necesarios para amortizarse por completo (en 6 ciclos solo se recuperarían 42,0 s).
\end{itemize}

\subsection{Contraste estadístico}
La prueba de Friedman confirmó diferencias significativas entre brazos en todas las métricas: $F_1$ macro ($\chi_F^2 = 25{,}32$, $p < 0{,}0001$), recall de ataque ($\chi_F^2 = 55{,}40$, $p < 0{,}0001$), accuracy ($\chi_F^2 = 20{,}28$, $p = 0{,}0001$), TFP ($\chi_F^2 = 10{,}67$, $p = 0{,}0137$) y TFD ($\chi_F^2 = 21{,}40$, $p = 0{,}0001$). La Tabla~\ref{tab:est} detalla los contrastes pareados post-hoc de Wilcoxon ajustados con Holm, incorporando el contraste crítico de tiempos frente al control aleatorio \emph{Random-k21}.

\begin{table}[ht]
\tabfont
\caption{Contrastes pareados post-hoc de Wilcoxon con corrección de Holm ($N = 30$)}
\label{tab:est}
\begin{tabular}{@{}llrrrc@{}}
\toprule
Comparación & Métrica & $p_{\text{raw}}$ & $p_{\text{Holm}}$ & Tamaño efecto ($r_b$) & Significativo ($\alpha=0{,}05$) \\
\midrule
PSO vs. Línea base & $F_1$ macro & 0,3050 & 1,0000 & $+0{,}019$ & No \\
PSO vs. Línea base & Recall (ataque) & $< 0{,}0001$ & $< 0{,}0010$ & $-1{,}000$ & \textbf{Sí (Línea base sup.)} \\
PSO vs. Línea base & Accuracy & 0,0570 & 0,3400 & $+0{,}465$ & No \\
PSO vs. Línea base & TFP & 0,0880 & 0,5264 & $-0{,}421$ & No \\
BA vs. Línea base  & $F_1$ macro & $< 0{,}0001$ & $< 0{,}0010$ & $-1{,}000$ & \textbf{Sí (Línea base sup.)} \\
BA vs. Línea base  & Recall (ataque) & $< 0{,}0001$ & $< 0{,}0010$ & $-1{,}000$ & \textbf{Sí (Línea base sup.)} \\
PSO vs. BA         & $F_1$ macro & $< 0{,}0001$ & $< 0{,}0010$ & $+1{,}000$ & \textbf{Sí (PSO sup.)} \\
PSO vs. BA         & Recall (ataque) & 0,0036 & 0,0215 & $-0{,}673$ & \textbf{Sí (BA sup.)} \\
\midrule
PSO vs. Línea base & $t_{\text{entrena}}$ & 0,00002 & $< 0{,}0010$ & $-0{,}884$ & \textbf{Sí (PSO más rápido)} \\
PSO vs. Línea base & $t_{\text{infer}}$ & 0,00015 & $< 0{,}0010$ & $-0{,}742$ & \textbf{Sí (PSO menor latencia)} \\
PSO vs. Random-k21 & $t_{\text{entrena}}$ & 0,0919 & 0,0919 & $-0{,}355$ & \textbf{No (Tiempos equivalentes)} \\
PSO vs. Random-k21 & $t_{\text{infer}}$ & 0,0262 & 0,0524 & $-0{,}462$ & No (tras Holm) \\
\bottomrule
\end{tabular}

\smallskip
\footnotesize\emph{Nota.} Convención de signo: $d = X_A - X_B$. En tiempos, $r_b < 0$ indica que el brazo $A$ toma menos tiempo que $B$.
\end{table}

\noindent\textbf{Hallazgo sobre el control aleatorio:}
La comparación directa de tiempos entre PSO ($16{,}91\text{ s}$) y \emph{Random-k21} ($16{,}73\text{ s}$) arrojó $p = 0{,}0919$ (no significativo). Esto demuestra que la ganancia del 36\,\% en tiempo de entrenamiento responde al menor número de atributos manipulados por el clasificador y no a una propiedad intrínseca de los subconjuntos seleccionados por PSO.

\subsection{Validación Leave-One-Capture-Out (LOO)}
En validación LOO, el $F_1$ macro decayó en todos los brazos debido al \emph{covariate shift} temporal: Línea base $= 0{,}486 \pm 0{,}031$; Random $= 0{,}480 \pm 0{,}025$; PSO $= 0{,}522 \pm 0{,}068$; BA $= 0{,}472 \pm 0{,}022$. PSO conservó mayor resiliencia en dos de las tres capturas evaluadas.

% =====================================================================
\section{Discusión}

\subsection{Comparación con la literatura}
Los hallazgos difieren marcadamente de las ganancias reportadas en bancos sintéticos. Mientras que \textcite{safa2024} y \textcite{ghanem2022} obtuvieron incrementos sistemáticos de precisión por encima del 98\,\% en NSL-KDD, en tráfico troncal real MAWI la selección mediante PSO apenas modificó el $F_1$ macro en $+0{,}86\,\%$ relativo sin significancia estadística ($p_{\text{Holm}} = 1{,}0000$). Los bancos sintéticos presentan fronteras de decisión artificiales de alta separabilidad, mientras que en el tráfico transpacífico crudo coexisten anomalías sutiles cuyo perfil se confunde con ráfagas benignas complejas, exigiendo la combinación sinérgica de múltiples atributos que los filtros metaheurísticos descartan.

\subsection{Relación con las hipótesis}
\textbf{H1 (Eficacia $F_1$ macro): NO respaldada.} PSO no difiere significativamente de la línea base ($p_{\text{Holm}} = 1{,}0000; r_b = +0{,}019$), y BA es significativamente inferior ($p_{\text{Holm}} < 0{,}001; r_b = -1{,}000$). Fundamentalmente, PSO provocó una caída estadísticamente significativa en el recall de la clase ataque de 6,39 puntos porcentuales (8,98\,\% relativo, de 0,7117 a 0,6478; $p_{\text{Holm}} < 0{,}001$), canjeando sensibilidad por especificidad y dejando inadvertidos a $13{,}127$ ataques adicionales por cada corrida de $900{,}000$ flujos.

\textbf{H2 (Falsas alarmas TFP): NO respaldada (formalmente).} Aunque PSO redujo la TFP de 0,1924 a 0,1516 ($-21{,}2\,\%$ relativo, evitando $28{,}275$ falsas alarmas por corrida), el ajuste de Holm arrojó $p_{\text{Holm}} = 0{,}5264$. No es posible rechazar la hipótesis nula con $\alpha = 0{,}05$ debido a la alta dispersión inter-semilla de PSO ($SD = 0{,}0588$ frente a $0{,}0003$ de la línea base).

\textbf{H3 (Eficiencia computacional): PARCIALMENTE respaldada.} Evaluando el umbral formal del 20\,\% fijado en la matriz de consistencia (así como el umbral más restrictivo del 25\,\%), el veredicto es idéntico: el tiempo de entrenamiento se redujo un 36,07\,\% ($p = 0{,}00002$), superando holgadamente ambos umbrales; sin embargo, la latencia de inferencia solo disminuyó 10,75\,\% ($p = 0{,}00015$), sin alcanzar la meta en ninguno de los dos casos. Además, la equivalencia estadística con el control aleatorio ($p = 0{,}0919$) ratifica que el ahorro proviene de la reducción dimensional y no del optimizador metaheurístico.

\textbf{H4 (Diferenciación cualitativa PSO vs. BA): PARCIALMENTE respaldada.} PSO y BA demostraron perfiles antagónicos de convergencia: PSO converge a soluciones compactas (18,5 atributos) con baja TFP, mientras que BA converge a subconjuntos más amplios (24,2 atributos) que preservan significativamente mayor recall de ataque ($p_{\text{Holm}} = 0{,}0215; r_b = -0{,}673$).

\subsection{Implicancias y limitaciones}
Reducir atributos no amplifica la detección en tráfico troncal y puede comprometer incidentes críticos al deteriorar el recall. Como limitaciones, las etiquetas proceden de MAWILab v1.1 (consenso no supervisado), se utilizaron capturas de 2021 del punto F y la búsqueda se acotó a submuestras de 10.000 flujos para viabilizar el tiempo de cómputo.

% ============================================================
\section{Conclusiones}
\begin{itemize}[leftmargin=1.5em,itemsep=2pt]
  \item \textbf{Hipótesis:} H1 no se respalda ($F_1$ no difiere, $p_{\text{Holm}} = 1{,}0000$; recall de ataque cae 6,39 pp / 8,98\,\%, $p_{\text{Holm}} < 0{,}001$); H2 no se respalda formalmente (TFP $-21{,}2\,\%$, $p_{\text{Holm}} = 0{,}5264$); H3 se respalda parcialmente (entrenamiento $-36{,}07\,\% \ge 20\,\%$, inferencia $-10{,}75\,\% < 20\,\%$); H4 se respalda parcialmente con perfiles de sensibilidad opuestos.
  \item \textbf{Relevancia y novedad:} Es una de las primeras evaluaciones controladas ($N = 30$ réplicas) en tráfico troncal transpacífico MAWI, contrastando críticamente con el optimismo metodológico de conjuntos sintéticos.
  \item \textbf{Aportes y avances:} Se constató empíricamente mediante un control aleatorio que el ahorro en tiempo de cómputo proviene de la reducción dimensional de atributos y no de la búsqueda bio-inspirada per se.
  \item \textbf{Aplicabilidad:} Aplicable en entornos donde el costo de búsqueda (33 ciclos de reentrenamiento para PSO y 7 ciclos para BA) sea amortizable frente a requerimientos estrictos de memoria.
  \item \textbf{Relación con la motivación:} El objetivo se cumplió al dimensionar de manera cuantitativa y transparente el comportamiento real de las metaheurísticas en IDS sobre enlaces troncales.
\end{itemize}

% =====================================================================
\section*{Dataset y código fuente}
\textbf{Dataset.} Los datos brutos provienen del repositorio público MAWI \parencite{mawiarchive} y etiquetas MAWILab \parencite{fontugne2010}. El conjunto consolidado en Parquet y los scripts de etiquetado están disponibles para investigación académica en el repositorio del proyecto.

\textbf{Código fuente.} La implementación completa en Python está disponible con instrucciones de replicabilidad en el repositorio: \url{https://github.com/Jhojan-Medina/thesis_ids_mawi}.

% =====================================================================
\printbibliography[title={Referencias}]

\end{document}
